\documentclass[a4paper,12pt]{article}
\usepackage[T2A]{fontenc}
\usepackage[utf8]{inputenc}
\usepackage[russian]{babel}
\usepackage{amsmath,amssymb,mathtools}
\usepackage{icomma}
\usepackage[a4paper,left=23mm,right=23mm,top=24mm,bottom=23mm,headheight=23pt,headsep=7mm]{geometry}
\usepackage{microtype}
\usepackage{tikz}
\usetikzlibrary{calc,arrows.meta,positioning,shapes.geometric}
\usepackage{tabularx,booktabs}
\usepackage{enumitem}
\usepackage{fancyhdr}
\usepackage{setspace}
\usepackage{titlesec}
\usepackage{needspace}
\usepackage{xcolor}
\usepackage[hidelinks]{hyperref}

\definecolor{accent}{HTML}{008080}
\definecolor{darkaccent}{HTML}{005A60}
\definecolor{textgray}{HTML}{30363A}
\definecolor{softgray}{HTML}{69727A}
\color{textgray}
\setstretch{1.22}
\setlength{\parindent}{0pt}
\setlength{\parskip}{5pt plus 1pt minus 1pt}
\setlength{\abovedisplayskip}{9pt plus 2pt minus 1pt}
\setlength{\belowdisplayskip}{10pt plus 2pt minus 1pt}
\setlength{\abovedisplayshortskip}{6pt}
\setlength{\belowdisplayshortskip}{6pt}
\setlist[itemize]{leftmargin=1.5em,itemsep=2pt,topsep=3pt,parsep=0pt}
\setlist[enumerate]{leftmargin=1.5em,itemsep=5pt,topsep=4pt,parsep=0pt}
\titleformat{\section}{\normalfont\color{darkaccent}\large}{\thesection.}{0.45em}{}[\vspace{4pt}\color{accent}\titlerule]
\titlespacing*{\section}{0pt}{22pt}{12pt}
\titleformat{\subsection}{\normalfont\color{darkaccent}\normalsize\bfseries}{}{0pt}{}
\titlespacing*{\subsection}{0pt}{11pt}{4pt}
\pagestyle{fancy}
\fancyhf{}
\fancyhead[L]{\footnotesize Алгебра и задачи на движение}
\fancyhead[R]{\footnotesize Ярослав Верещагин}
\fancyfoot[C]{\footnotesize\thepage}
\renewcommand{\headrulewidth}{0.4pt}
\renewcommand{\headrule}{\hbox to\headwidth{\color{accent}\leaders\hrule height \headrulewidth\hfill}}
\renewcommand{\footrulewidth}{0pt}
\newcommand{\method}[1]{\par\medskip\textcolor{darkaccent}{\textbf{#1}}\par\smallskip}
\newcommand{\important}[1]{\textcolor{darkaccent}{#1}}

\begin{document}
\thispagestyle{empty}
\begin{center}
\vspace*{18mm}
{\large\color{darkaccent} МАТЕМАТИКА\par}
\vspace{16mm}
{\LARGE Чек-лист\par}
\vspace{13mm}
{\LARGE\color{accent}\begin{tabular}{c}
Алгебраические преобразования\\[3mm]
и задачи на движение
\end{tabular}\par}
\vspace{13mm}
{\normalsize Подготовка к ЕГЭ по профильной математике\par}
\vspace{24mm}
{\normalsize Ярослав Верещагин\par}
\vspace{11mm}
{\normalsize 23 ноября 2026 г.\par}
\vfill
{\footnotesize Лёвин Артём Александрович\par}
{\footnotesize эксклюзивно для Ярослава Верещагина\par}
\vspace{10mm}
\begin{tikzpicture}[x=1cm,y=1cm]
\draw[color=accent,line width=.75pt]
 (0,0) .. controls (1.25,0.58) and (2.15,-0.32) .. (3.55,0.04)
 .. controls (4.55,0.36) and (5.17,0.53) .. (6.35,0.10)
 .. controls (7.25,-0.18) and (8.42,-0.24) .. (9.3,0.12);
\fill[accent] (4.82,0.35) circle (0.022);
\end{tikzpicture}
\end{center}
\clearpage
\setcounter{page}{1}

\section{Раскрытие скобок без потери знаков}

Если перемножаются две суммы, \important{каждое слагаемое первой скобки} умножается на \important{каждое слагаемое второй}. Четыре произведения сначала удобно выписать отдельно:
\[
(a+b)(c+d)=ac+ad+bc+bd.
\]

\method{Пример 1. Перемножение двучленов}
\begin{align*}
(3x+5)(1-2x)
 &=3x\cdot1+3x\cdot(-2x)+5\cdot1+5\cdot(-2x)\\
 &=3x-6x^2+5-10x\\
 &=-6x^2-7x+5.
\end{align*}

Обратите внимание: при умножении $3x\cdot(-2x)$ получается $-6x^2$, а не $-6x$. Переменные тоже перемножаются.

\method{Пример 2. Два отрицательных множителя}
\begin{align*}
(-4x-1)(10x-7)
 &=-40x^2+28x-10x+7\\
 &=-40x^2+18x+7.
\end{align*}

\textbf{Проверка в два прохода:} сначала соберите подобные слагаемые, затем заново посмотрите только на знаки всех четырёх произведений.

\section{Таблица для задачи на движение}

Три связанные величины: \important{расстояние $s$}, \important{скорость $v$} и \important{время $t$}.
\[
 s=vt,\qquad t=\frac{s}{v},\qquad v=\frac{s}{t}.
\]

Удобно сразу заполнить таблицу, а затем перевести условие о разности или сумме времён в уравнение.

\method{Пример 3. Движение туда и обратно}

Расстояние между пунктами $A$ и $B$ равно $28$ км. В обратном направлении скорость на $9$ км/ч выше, а время движения на $1$ ч меньше. Найдите скорость в обратном направлении.

Пусть $x$ км/ч --- скорость из $A$ в $B$. Тогда $x>0$ и:
\begin{center}
\renewcommand{\arraystretch}{1.28}
\begin{tabular}{@{}lccc@{}}
\toprule
Направление & Путь, км & Скорость, км/ч & Время, ч\\
\midrule
Из $A$ в $B$ & $28$ & $x$ & $\dfrac{28}{x}$\\
Из $B$ в $A$ & $28$ & $x+9$ & $\dfrac{28}{x+9}$\\
\bottomrule
\end{tabular}
\end{center}

Первое время больше второго на час:
\[
\frac{28}{x}-\frac{28}{x+9}=1.
\]

\section{Как получить и решить квадратное уравнение}

Продолжим пример с движением. Для $x>0$ оба знаменателя положительны, поэтому можно умножить уравнение на $x(x+9)$:
\begin{align*}
28(x+9)-28x&=x(x+9),\\
252&=x^2+9x,\\
x^2+9x-252&=0.
\end{align*}

Для уравнения $ax^2+bx+c=0$, $a\ne0$, используем:
\[
D=b^2-4ac,\qquad x_{1,2}=\frac{-b\pm\sqrt D}{2a}\quad(D\ge0).
\]

Здесь
\[
D=9^2-4\cdot1\cdot(-252)=1089=33^2,
\]
\[
x_1=\frac{-9+33}{2}=12,\qquad
x_2=\frac{-9-33}{2}=-21.
\]

Отрицательное значение $-21$ не подходит по смыслу задачи: скорость не может быть отрицательной. Найденное $x=12$ --- скорость \emph{туда}; в ответ нужен обратный путь:
\[
\boxed{v_{B\to A}=x+9=21\text{ км/ч}.}
\]

Проверка: $28/12-28/21=1$ ч.

\Needspace{14\baselineskip}
\section{Как извлекать корни из больших квадратов}

\method{Способ 1. Границы и последняя цифра}

Для $\sqrt{1089}$ имеем $30^2<1089<40^2$: целый корень, если он существует, находится между $30$ и $40$. На конце числа стоит $9$, значит последняя цифра корня может быть $3$ или $7$. Проверяем кандидатов:
\[
33^2=1089,\qquad\boxed{\sqrt{1089}=33}.
\]

Другой пример: $70^2<5184<80^2$, а последняя цифра $4$ даёт два кандидата --- $72$ и $78$:
\[
72^2=5184,\qquad\boxed{\sqrt{5184}=72}.
\]

\textbf{Важно:} границы и последняя цифра позволяют найти \emph{кандидатов}, но окончательный ответ подтверждаем возведением в квадрат.

\method{Способ 2. Выделение полного квадрата}
\[
\sqrt{19600}=\sqrt{196\cdot100}=14\cdot10=\boxed{140}.
\]

\method{Способ 3. Число около знакомого квадрата}

В уравнении $x^2+1005x+3006=0$ дискриминант равен
\[
D=1005^2-4\cdot3006=998001=999^2.
\]

Действительно, $(1000-1)^2=1000000-2000+1=998001$. Тогда корни:
\[
 x_1=\frac{-1005+999}{2}=-3,\qquad
 x_2=\frac{-1005-999}{2}=-1002.
\]

\section{Средняя скорость: общий принцип}

Средняя скорость --- \important{весь пройденный путь}, делённый на \important{всё затраченное время}:
\[
\boxed{v_{\text{ср}}=\frac{s_1+s_2+\cdots+s_n}{t_1+t_2+\cdots+t_n}}.
\]

Обычное среднее арифметическое скоростей подходит не всегда. Нужно увидеть, что по условию разделено на равные части: \emph{время} или \emph{расстояние}.

\subsection{Случай 1. Равные промежутки времени}

Половину времени автомобиль ехал со скоростью $84$ км/ч, оставшуюся половину --- $56$ км/ч. Обозначим всё время $t>0$.
\begin{center}
\renewcommand{\arraystretch}{1.22}
\begin{tabular}{@{}lccc@{}}
\toprule
Часть поездки & Скорость & Время & Путь\\
\midrule
Первая & $84$ & $t/2$ & $42t$\\
Вторая & $56$ & $t/2$ & $28t$\\
\bottomrule
\end{tabular}
\end{center}
\[
 v_{\text{ср}}=\frac{42t+28t}{t/2+t/2}
 =\frac{70t}{t}=\boxed{70\text{ км/ч}}.
\]

Для двух одинаковых временных отрезков можно сразу взять
\[
\boxed{v_{\text{ср}}=\frac{v_1+v_2}{2}}.
\]

\subsection{Случай 2. Равные участки пути}

Первую половину трассы автомобиль ехал со скоростью $72$ км/ч, вторую --- $108$ км/ч. Обозначим всю длину трассы $s>0$.
\begin{center}
\renewcommand{\arraystretch}{1.22}
\begin{tabular}{@{}lccc@{}}
\toprule
Часть поездки & Скорость & Путь & Время\\
\midrule
Первая & $72$ & $s/2$ & $s/144$\\
Вторая & $108$ & $s/2$ & $s/216$\\
\bottomrule
\end{tabular}
\end{center}
\[
 v_{\text{ср}}=\frac{s}{\dfrac{s}{144}+\dfrac{s}{216}}
 =\frac{144\cdot216}{144+216}
 =\boxed{86{,}4\text{ км/ч}}.
\]

Для \emph{двух равных расстояний} получается другая формула:
\[
\boxed{v_{\text{ср}}=\frac{2v_1v_2}{v_1+v_2}}.
\]

\section{Работа со сложными дробями}

Делить на дробь означает умножать на обратную. Если знаменатель --- обычное число, мысленно представьте его дробью с единицей:
\[
\frac{\dfrac ab}{\dfrac cd}=\frac{ad}{bc},\qquad
\frac{\dfrac ab}{c}=\frac{a}{bc},\qquad b,c,d\ne0.
\]

\clearpage
\thispagestyle{fancy}
\section*{Алгоритм: как выбрать формулу средней скорости}
\addcontentsline{toc}{section}{Алгоритм средней скорости}

\vspace{8mm}
\begin{center}
\begin{tikzpicture}[
  >={Latex[length=2.7mm]},
  line width=.85pt,
  font=\normalsize,
  start/.style={ellipse,draw=accent,thick,align=center,minimum width=7.8cm,minimum height=1.2cm},
  action/.style={rounded corners=2mm,draw=accent,thick,align=center,text width=6.25cm,inner sep=7pt,minimum height=1.45cm},
  choice/.style={diamond,aspect=3.0,draw=accent,thick,align=center,text width=4.2cm,inner sep=3pt},
  finish/.style={ellipse,draw=accent,thick,align=center,minimum width=9cm,minimum height=1.25cm},
  arrow/.style={->,draw=darkaccent,thick}
]
\node[start] (begin) at (0,0) {Начните с общего правила\\$v_{\text{ср}}=S_{\text{общ}}/t_{\text{общ}}$};
\node[action] (table) at (0,-2.20) {Запишите для каждого участка\\скорость, время и расстояние};
\node[choice] (choice) at (0,-4.85) {Что разделено на две\\равные части?};
\node[action,text width=5.35cm] (time) at (-4.00,-8.10) {\textbf{Время}\\[2pt] $v_{\text{ср}}=\dfrac{v_1+v_2}{2}$};
\node[action,text width=5.35cm] (dist) at (4.00,-8.10) {\textbf{Расстояние}\\[2pt] $v_{\text{ср}}=\dfrac{2v_1v_2}{v_1+v_2}$};
\node[action,text width=10.4cm] (check) at (0,-11.05) {Убедитесь, что результат между $v_1$ и $v_2$;\\проверьте единицы измерения};
\node[finish] (end) at (0,-13.20) {Запишите ответ в км/ч};
\draw[arrow] (begin) -- (table);
\draw[arrow] (table) -- (choice);
\draw[arrow] (choice.west) -| node[pos=.32,above,sloped,font=\small]{время} (time.north);
\draw[arrow] (choice.east) -| node[pos=.32,above,sloped,font=\small]{путь} (dist.north);
\coordinate (merge) at (0,-9.67);
\draw[draw=darkaccent,thick] (time.south) |- (merge);
\draw[draw=darkaccent,thick] (dist.south) |- (merge);
\draw[arrow] (merge) -- (check.north);
\draw[arrow] (check) -- (end);
\end{tikzpicture}
\end{center}
\clearpage

\section{Что нужно проверять самостоятельно}

\begin{itemize}[leftmargin=1.8em,itemsep=7pt]
\item \textbf{Скобки.} Все ли четыре произведения выписаны? Не потерян ли $x$ или $x^2$?
\item \textbf{Знаки.} Действительно ли произведение двух отрицательных чисел положительно?
\item \textbf{Дискриминант.} Учитывается ли знак коэффициента $c$ в формуле $D=b^2-4ac$?
\item \textbf{Корень.} Проверен ли найденный целый корень возведением в квадрат?
\item \textbf{Движение.} Не перепутаны ли путь, скорость и время? Положительна ли скорость?
\item \textbf{Ответ.} Спрашивают скорость туда или обратно, среднюю скорость или время?
\item \textbf{Средняя скорость.} Равны части по времени или по пути? Не усреднены ли скорости там, где это неверно?
\item \textbf{Дроби.} Отмечены ли ненулевые знаменатели и обосновано ли сокращение?
\end{itemize}

\method{Краткая памятка}
\begin{center}
\renewcommand{\arraystretch}{1.35}
\begin{tabularx}{\linewidth}{@{}>{\raggedright\arraybackslash}p{.34\linewidth} X@{}}
\toprule
Сигнал в условии & Полезное действие\\
\midrule
Произведение скобок & Выпишите четыре попарных произведения; затем приведите подобные.\\
Большой дискриминант & Найдите близкий квадрат или выделите полный квадрат в произведении.\\
Разность времён & Составьте таблицу $s,v,t$ и уравнение по условию.\\
Половина \emph{времени} & Для двух равных промежутков используйте среднее арифметическое.\\
Половина \emph{пути} & Складывайте времена; для двух равных участков используйте $2v_1v_2/(v_1+v_2)$.\\
Дробь в знаменателе & Замените деление умножением на обратную дробь.\\
\bottomrule
\end{tabularx}
\end{center}

\clearpage
\section*{Тренировочный блок}
\addcontentsline{toc}{section}{Тренировочный блок}
\begingroup
\setstretch{1.05}
\setlength{\abovedisplayskip}{5pt}
\setlength{\belowdisplayskip}{5pt}
\textcolor{softgray}{Путь А · 10 задач · от базового уровня к повышенному. Решения запишите по шагам; для текстовых задач составьте таблицу величин.}

\subsection*{Базовый уровень}
\begin{enumerate}[label=\arabic*.,leftmargin=1.8em,itemsep=4pt]
\item Раскройте скобки и приведите подобные:
\[
(2x+5)(1-3x).
\]
\item Раскройте скобки, внимательно проверяя знаки:
\[
(-3x-2)(5x-4).
\]
\item Вычислите сложную дробь:
\[
\frac{\dfrac34}{\dfrac56}.
\]
\end{enumerate}

\subsection*{Средний уровень}
\begin{enumerate}[label=\arabic*.,leftmargin=1.8em,start=4,itemsep=4pt]
\item Вычислите $\sqrt{7056}$, предварительно ограничив корень соседними десятками и проверив последнюю цифру.
\item Вычислите $\sqrt{32400}$, выделив полный квадрат в произведении.
\item Решите квадратное уравнение:
\[
x^2+11x-312=0.
\]
\item Решите уравнение, используя удобный способ извлечения корня из дискриминанта:
\[
x^2+1003x+2002=0.
\]
\end{enumerate}

\subsection*{Уровень выше среднего}
\begin{enumerate}[label=\arabic*.,leftmargin=1.8em,start=8,itemsep=6pt]
\item Первую половину \emph{времени} автомобиль двигался со скоростью $48$ км/ч, вторую половину --- $72$ км/ч. Найдите среднюю скорость за всю поездку.
\item Первую половину \emph{пути} автомобиль проехал со скоростью $60$ км/ч, вторую --- $90$ км/ч. Найдите среднюю скорость за всю поездку.
\item Между пунктами $A$ и $B$ --- $36$ км. Обратно автомобиль едет на $6$ км/ч быстрее, затрачивая на $1$ ч меньше. Найдите скорость на \emph{обратном} пути. Составьте таблицу и объясните, почему отрицательный корень не подходит.
\end{enumerate}
\endgroup

\clearpage
\section*{Ответы к тренировочному блоку}
\addcontentsline{toc}{section}{Ответы}
\vspace{3mm}
\begin{center}
\renewcommand{\arraystretch}{1.55}
\begin{tabularx}{\linewidth}{@{}cX@{}}
\toprule
№ & Ответ\\
\midrule
1 & $-6x^2-13x+5$\\
2 & $-15x^2+2x+8$\\
3 & $\dfrac9{10}$\\
4 & $84$\\
5 & $180$\\
6 & $x_1=13$, $x_2=-24$\\
7 & $x_1=-2$, $x_2=-1001$\\
8 & $60$ км/ч\\
9 & $72$ км/ч\\
10 & $18$ км/ч (скорость из $A$ в $B$ --- $12$ км/ч)\\
\bottomrule
\end{tabularx}
\end{center}

\vspace{16mm}
\textcolor{darkaccent}{\textbf{Самопроверка после решения}}
\begin{itemize}[leftmargin=1.7em,itemsep=5pt]
\item Для заданий 1--2 раскройте произведения повторно, отдельно отслеживая знаки.
\item Для заданий 4--5 возведите полученные целые числа в квадрат.
\item Для заданий 6--7 подставьте оба корня в исходное уравнение.
\item Для заданий 8--10 проверьте, что средняя скорость лежит между скоростями участков, а найденная скорость удовлетворяет условию о разности времён.
\end{itemize}

\vfill
\begin{center}
{\footnotesize\color{softgray}Лёвин Артём Александрович эксклюзивно для Ярослава Верещагина}
\end{center}
\end{document}
